% =====================================================================
% Module 05 — Composite Operator Algebra
% Academic-Galactic v1.2
% =====================================================================

\section{Composite Operator Algebra}
\label{sec:CompositeOperatorAlgebra}

\subsection{Overview}
The composite operator algebra defines the operator backbone of the Soft
80-dimensional manifold $\mathcal{M}^{80}$. It integrates:

\begin{itemize}
    \item spectral operator $\Phi_5$ (Module 03),
    \item curvature envelope operator $\Phi_7$ (Module 02),
    \item holonomy operator $\Phi_8$ (Module 02),
    \item topology-modulation drift operator $\Phi_8^t$ (Module 02),
    \item completion operator $\Phi_9$ (Module 08),
    \item domain operator $\mathcal{D}$ (Module 01),
    \item altitude operator $\mathcal{A}$ (Module 01),
    \item compatibility operator $\mathcal{C}$ (Module 06),
    \item reprojection operator $\mathcal{R}$ (Module 07).
\end{itemize}

Operator-algebraic foundations follow state-space geometry
\citeClassA{AlfsenSchultz2003} and modern categorical-homotopical operator
algebra research \citeClassD{EmmersonBuchanan2025}.

\subsection{Operator Set}
The composite operator algebra is defined as:
\[
\mathcal{A}_{\mathrm{comp}}
=
\left\{
\Phi_5, \Phi_7, \Phi_8, \Phi_8^t, \Phi_9,
\mathcal{D}, \mathcal{A}, \mathcal{C}, \mathcal{R}
\right\}.
\]

Each operator acts on filtration layers $\mathcal{F}_k$ and their associated
fiber and spectral bundles.

\subsection{Operator Composition}
Operators compose according to:
\[
\mathcal{O}_2 \circ \mathcal{O}_1 : \mathcal{F}_k \to \mathcal{F}_{k'},
\]
where altitude transitions are governed by $\mathcal{A}$ and domain transitions
by $\mathcal{D}$.

Composition is associative:
\[
(\mathcal{O}_3 \circ \mathcal{O}_2) \circ \mathcal{O}_1
=
\mathcal{O}_3 \circ (\mathcal{O}_2 \circ \mathcal{O}_1).
\]

\subsection{Lie Bracket Structure}
The operator algebra admits a Lie bracket:
\[
[\mathcal{O}_1, \mathcal{O}_2]
=
\mathcal{O}_1 \circ \mathcal{O}_2
-
\mathcal{O}_2 \circ \mathcal{O}_1.
\]

Important relations include:
\[
[\Phi_7, \Phi_8] = 0
\quad\text{(holonomy generated by curvature)},
\]
\[
[\Phi_8, \Phi_8^t] \neq 0
\quad\text{(drift induces topology modulation)},
\]
\[
[\Phi_5, \mathcal{C}] = 0
\quad\text{(spectral–dual alignment)}.
\]

\subsection{Filtration Preservation}
Operators preserve filtration structure:
\[
\mathcal{O}(\mathcal{F}_k) \subseteq \mathcal{F}_{k+1}.
\]

This follows from accessible category structure
\citeClassA{Adamek1994}
and topos-theoretic stratification
\citeClassC{BarrWells2000,Bell1982}.

\subsection{Curvature–Holonomy Interaction}
Curvature envelopes generate holonomy:
\[
\Phi_8(\gamma)
=
\exp\left( \int_\gamma \Phi_7 \right),
\]
consistent with classical holonomy theory
\citeClassA{Atiyah1956}
and modern holonomy-curvature research
\citeClassA{BlackholeHolonomy2026,GutierrezUle2026}.

\subsection{Topology-Modulation Drift}
The drift operator $\Phi_8^t$ modulates topology:
\[
\Phi_8^t : \mathcal{F}_k \to \mathcal{F}_k.
\]

It interacts with collapse surfaces (Module 04) and satisfies:
\[
[\Phi_8^t, \Phi_7] \neq 0,
\]
indicating curvature-driven topology change.

\subsection{Spectral–Dual Compatibility}
Dual compatibility requires:
\[
\mathcal{C}(\Phi_5(p)) = \Phi_5(\mathcal{C}(p)),
\]
ensuring cotangent spectral consistency (Module 06).

\subsection{Reprojection Idempotence}
The reprojection operator $\mathcal{R}$ satisfies:
\[
\mathcal{R}^2 = \mathcal{R},
\]
ensuring geometric consistency restoration (Module 07).

\subsection{Completion Operator}
The completion operator $\Phi_9$ extracts invariant structure:
\[
\Phi_9 : \mathcal{F}_k \to \mathcal{F}_k,
\]
and satisfies:
\[
[\Phi_9, \Phi_5] = 0,
\qquad
[\Phi_9, \Phi_7] = 0.
\]

\subsection{Summary}
The composite operator algebra defines the operator backbone of the Soft 80D
manifold, integrating spectral, curvature, holonomy, drift, dual, domain,
altitude, reprojection, and completion operators. It provides the algebraic
structure required for the A7 operator altitude and ensures compatibility with
filtration, fiber bundle geometry, spectral continuity, collapse-surface
stratification, and completion conditions.

% =====================================================================
% End Module 05
% =====================================================================
